\documentclass[openany,fontset=fandol]{ctexbook}
\usepackage{amsmath,amssymb,amsthm}
\usepackage[unicode,hidelinks]{hyperref}
\setlength{\emergencystretch}{2em}
\raggedbottom
\DeclareMathOperator{\Spec}{Spec}
\begin{document}
\title{代数几何基础第五卷预备笔记\\法文覆盖审计：有限 étale 多截面}
\author{第99批\quad 法、英、中文三方校订}
\date{2026年11月24日}
\maketitle

\section*{范围与接续}
连续法校至物理29页末，第7节结束。对应英物理20页命题7.4、21页开头备注7.5；另核英46页关于余维 $n$ 线性截面的作者补记。本批未改1158页合订本；下一批从法物理30页第8节开始。

\section*{7.4\quad 原稿的命题与证明设想}
\textbf{命题（手稿所述）.}
设 $g:X\to S$ 固有且平坦，取 $s\in S$，令 $n=\dim X_s$。假设 $X_s$ 上不对 $k(s)$ 可分的点所成集合，其维数\textbf{严格小于} $n$（例如 $X_s$ 可分）。那么存在 $s$ 的开邻域 $U$ 和有限、étale、满射的 $S'\to U$，使 $X\times_S S'\to S'$ 有截面。若 $k(s)$ 无限，原稿还说可取 $S'$ 为 $X_U$ 的闭子概形。\footnote{法29清楚印为“$<n$”，英20误作“$\le n$”；下文给出反例。断言及新增旁注均为手稿材料，不能略去下面指出的适用范围。}

\textbf{原稿的证明路线.}
先设 $k(s)$ 无限。对 $n$ 归纳：$n=0$ 时，原稿称缩小 $U$ 后 $X_U\to U$ 本身即有限、étale、满射。$n>0$ 时，作者提议对“可分”应用7.3，用一个超平面截面 $Y$ 替代 $X$，使维数由 $n$ 降为 $n-1$ 并保留所需条件。若未假设 $k(s)$ 无限，则先作适当的 étale 换底。

这段归纳需要有可用的射影嵌入与控维条件；仅由固有和平坦并不能得到这种超平面截面。因此保留为原稿的证明路线，不据此认定一般固有情形已经证明。下文另给出有明确附加假设的版本。

\section*{7.5\quad 关于射影性的备注}
作者指出，射影且相对可分的族，在所讨论的非空纤维附近可有 étale 多截面。作者接着提到，若只固有、光滑而不射影，同样结论可能失败；并提出可从非有限的 étale 态射寻找拟射影情形下的反例。\footnote{法29这一段只写明寻找反例的方向；页末手记的末尾读法未定。没有给出一个完整的固有、光滑而非射影的构造。非有限的 étale 态射本身不足以说明无有限 étale 多截面。若仍同时要求固有，则拟射影态射已经是射影的。}

\section*{译校补证}
\paragraph{英译放宽为 $\le n$ 会出错.}
取 $k=\mathbf F_p(t)$，令 $X=\Spec(k[u]/(u^p-t))$、$S=\Spec k$。
由 $t$-进赋值 $v_t(t)=1$，$t$ 不是 $k$ 的 $p$ 次幂；该代数是次数 $p$ 的纯不可分域扩张。因此 $X\to S$ 有限、平坦且固有。$n=0$，唯一点不具几何约化性，坏集也是零维，满足英译的 $\le n$。然而若 $S'\to S$ 非空且有限 étale，它是有限可分域扩张的谱之有限不交并，其中任何一个域都不能含 $t$ 的 $p$ 次根。因此换底后无截面。反例不适用于法稿严格不等号。

\paragraph{零维基例与较弱的局部结论.}
若 $g$ 有限表示且 $X_s$ 非空零维，严格 $<0$ 使坏集为空。于是 $X_s$ 有限 étale。由平坦和有限表示，纤维上的点都属于 $g$ 的 étale 轨迹；利用固有性，把其补集的闭像从底去掉。缩小后 $g$ 固有且 étale，因而固有且拟有限，所以有限；再利用 étale 态射是开的，把底限制到其包含 $s$ 的像即可得到满射。一般维数时，如果 $X_s$ 中存在一个 $g$-光滑点，则有经适当 étale 邻域的局部截面（\href{https://stacks.math.columbia.edu/tag/055U}{Stacks 37.38.6}）；这条结论并不把该邻域变成底开集的有限 étale 覆盖。

\paragraph{一个可核验的射影版本.}
假定 $g$ 射影、平坦、有限表示，$X_s$ 非空且纯维 $n$，并假定 $X_s$ 的非光滑轨迹维数小于 $n$。选取嵌入 $X\hookrightarrow\mathbf P^N_S$。在几何纤维上，余维 $n$ 的线性子空间中，与非光滑轨迹相交者构成 Grassmann 丛参数空间的真闭子集：对固定点，包含该点施加 $n$ 个条件，故上述轨迹的维数不足以覆盖参数空间。在光滑点上，包含该点的子空间若还不横截切空间，再多至少一个独立的闭条件。将该关联坏轨迹投到参数空间，其闭包维数仍严格小于参数空间。因而存在非空开集，其任意参数给出非空、有限、几何约化的横截；非空由射影维数不等式保证。这也解释了英46页作者提出直接取余维 $n$ 的截面而不用逐次切割。

若 $k(s)$ 无限，利用 Grassmann 丛的仿射大图可在该开集中选 $k(s)$-点，并在 $s$ 的开邻域提升矩阵坐标。所得闭截面在纤维上的各点处，由底上的光滑性及方程微分线性无关，对底 étale。该截面相对底固有，所以可缩小底去除非 étale 点的闭像。由固有且拟有限得有限，由非空纤维与开映射得缩小后的满射。取此截面自身为 $S'$，闭嵌入给出所需的换底截面。若 $k(s)$ 有限，可先取好参数开集中的一个闭点；其剩余域有限可分。按7.1作纤维只有该点的有限 étale 换底，提升其参数。若需在新底缩邻域，可先从原底删去新底坏闭集在有限态射下的闭像；最终有限 étale 满射可与第一次换底复合。

\clearpage
\paragraph{拟射影但无此类多截面的具体反例.}
令 $A=\mathbf Q[t]_{(t)}$，$S=\Spec A$，令
\[
 D=\Spec\bigl(A[z]/(z^2-(1+t))\bigr).
\]
$1+t$ 是单位，$2z$ 在该代数中可逆，故 $D\to S$ 有限 étale。$1+t$ 在 $\mathbf Q(t)$ 不是平方，而该代数是自由 $A$-模，故它嵌入其泛纤维的域，因而 $D$ 整且连通。闭纤维有两点 $z=\pm1$；去掉其中一点得 $X\subset D$。于是 $X\to S$ 是满射、拟射影、étale 的，但不固有。

若存在有限 étale 满射 $T\to S$ 和一个 $S$-态射 $T\to X$，则 $T\to D$ 有限（两者都仿射有限于 $S$），也 étale（\href{https://stacks.math.columbia.edu/tag/0HB1}{Stacks 29.37.18}）。其像非空且在连通的 $D$ 中既开又闭，只能是整个 $D$，与被删点不在像中矛盾。$S$ 的任一含闭点的开邻域已是整个 $S$，故缩小邻域也无法补救。这是对“以拟射影替代射影”并丢掉固有性的反例，不是对一般固有、光滑非射影情形的反例。
\end{document}

